
\documentclass[12pt]{article}
\usepackage{amssymb}
\usepackage{amsmath}

\setlength{\topmargin}{-1.0in}
\setlength{\textheight}{9.25in}
\setlength{\oddsidemargin}{0.0in}
\setlength{\evensidemargin}{0.0in}
\setlength{\textwidth}{6.5in}
\def\labelenumi{\arabic{enumi}.}
\def\theenumi{\arabic{enumi}}
\def\labelenumii{(\alph{enumii})}
\def\theenumii{\alph{enumii}}
\def\p@enumii{\theenumi.}
\def\labelenumiii{\arabic{enumiii}.}
\def\theenumiii{\arabic{enumiii}}
\def\p@enumiii{(\theenumi)(\theenumii)}
\def\labelenumiv{\arabic{enumiv}.}
\def\theenumiv{\arabic{enumiv}}
\def\p@enumiv{\p@enumiii.\theenumiii}
\pagestyle{plain}
\setcounter{secnumdepth}{0}
\parindent=0pt
\begin{document}


\begin{center}
\textbf{Econ 326}

\textbf{Assignment 6}
\end{center}

\begin{enumerate}


\item An econometrician used R to regress the dependent variable \verb!liver! on the independent variable \verb!alcohol!, where \verb!liver! is the number of liver disease deaths per 100,000 people in a country, and \verb!alcohol! is consumption of alcohol in liters per capita in a country. The following output was obtained:
\small
\begin{verbatim}
> model <- lm(liver ~ alcohol, data = mydata)
> summary(model)
...
Coefficients:
            Estimate Std. Error t value Pr(>|t|)
(Intercept)  10.8548     2.8024   3.874 0.001030 **
alcohol       3.5864     0.7541       A        B
---
Signif. codes:  0 '***' 0.001 '**' 0.01 '*' 0.05 '.' 0.1 ' ' 1

Residual standard error: 8.29 on 19 degrees of freedom


> confint(model)
                2.5 %   97.5 %
(Intercept) 4.989313 16.72033
alcohol            C        D
\end{verbatim}
\normalsize
\begin{enumerate}
\item Several entries in the output were replaced with letters. Find A--D. Show your work.
\item Test at 5\% significance level that the coefficient of \verb!alcohol! is 5 (against the alternative that it is different from 5). \label{it:test}
\item Test the same hypothesis as in part (b) at 10\% significance level.
\end{enumerate}

\item An econometrician used R to regress the dependent variable \verb!ln_heart! on the independent variable \verb!ln_alcohol!, where \verb!ln_heart! is the natural logarithm of the number of heart disease deaths per 100,000 people in a country, and \verb!ln_alcohol! is the natural logarithm of consumption of alcohol in liters per capita in a country. The following output was obtained:
\small
\begin{verbatim}
> model <- lm(ln_heart ~ ln_alcohol, data = mydata)
> summary(model)
...
Coefficients:
             Estimate Std. Error t value Pr(>|t|)
(Intercept)  5.361366   0.132482  40.467   <2e-16 ***
ln_alcohol          A          B       C        D
...

> confint(model)
                  2.5 %     97.5 %
(Intercept)   5.084078   5.638654
ln_alcohol   -0.611437  -0.095217
\end{verbatim}
\normalsize
\begin{enumerate}
\item  Several entries in the output were replaced with letters. Find A--D if the number of observations was 21.
 Provide careful explanations for your answers.
\item  What is the interpretation of the coefficient on \verb!ln_alcohol!? Explain.
\item  Test at 5\% significance level $H_0: \beta \geq 0$ against $H_1: \beta<0$, where $\beta$ is the coefficient on \verb!ln_alcohol!.
\end{enumerate}

\item An econometrician used R to regress \verb!y! on \verb!x!. The
following output was obtained:
\small
\begin{verbatim}
> model <- lm(y ~ x, data = mydata)
> summary(model)
...
Coefficients:
            Estimate Std. Error t value Pr(>|t|)
(Intercept) 3.109574   0.182132  17.073   <2e-16 ***
x           0.304925   0.177374       A        B
---
Signif. codes:  0 '***' 0.001 '**' 0.01 '*' 0.05 '.' 0.1 ' ' 1

Residual standard error: 0.91 on 23 degrees of freedom
Multiple R-squared:  0.1139, Adjusted R-squared:  0.0753
F-statistic: 2.96 on 1 and 23 DF,  p-value: 0.09900

> confint(model)
                 2.5 %    97.5 %
(Intercept)  2.732806  3.486343
x                   C         D
\end{verbatim}
\normalsize
\begin{enumerate}
\item Find A--D. Provide careful explanations for your answers.
\item Let $\beta$ denote the coefficient on \verb!x!. Test at
5\% significance level $H_0:\beta=0$ against $H_1:\beta
\ne 0$.
\item Test at 5\% significance level $H_0:\beta \leq 0$ against $H_1:\beta
> 0$.
\item Explain the difference between the outcomes in parts (b)
and (c).
\end{enumerate}

\item Consider testing the null hypothesis $H_0\colon \beta = 1$ at the 5\% significance level using two different tests:
\begin{itemize}
\item \textbf{Two-sided test:} $H_0\colon \beta = 1$ vs.\ $H_1\colon \beta \ne 1$
\item \textbf{One-sided test:} $H_0\colon \beta = 1$ vs.\ $H_1\colon \beta > 1$
\end{itemize}
Suppose that $\text{Var}(\hat{\beta}) = 4$ (known), so $\sigma_{\hat{\beta}} = 2$. Since the variance is known, use the standard normal distribution.
\begin{enumerate}
\item Find the rejection regions for both tests in terms of $\hat{\beta}$.
\item Compute the power of the two-sided test for each of the following true values of $\beta$: $\beta^* \in \{-3,\; 0,\; 2,\; 5\}$.
\item Compute the power of the one-sided test for each of the same true values: $\beta^* \in \{-3,\; 0,\; 2,\; 5\}$.
\item Compare the power of the two tests. When does the one-sided test have higher power? When does it have lower power? Explain.
\end{enumerate}

\item Consider a two-sided test of $H_0\colon \beta = 0$ against
$H_1\colon \beta \ne 0$ at the 5\% significance level. Assume the
variance of $\hat{\beta}$ is known, so the standard normal
distribution applies.

When $\beta^* > 0$, the power of the two-sided test is approximately
$$
\text{Power} \approx 1 - \Phi\!\left(z_{0.975} -
\frac{\beta^*}{\sigma_{\hat{\beta}}}\right),
$$
where $\Phi$ is the standard normal CDF and $z_{0.975} = 1.96$.
This approximation ignores the small probability of rejecting
in the left tail, which is negligible when $\beta^*>0$.

Suppose a researcher wants the test to have at least 80\% power.
The smallest value of $|\beta^*|$ for which this holds is called
the \emph{minimum detectable effect} (MDE): any true effect
smaller than the MDE is unlikely to be detected by the test.

\medskip
\textit{Hint:} $\Phi^{-1}(0.20) = -0.84$, equivalently $z_{0.80} = 0.84$.

\begin{enumerate}
\item Find the MDE when $\sigma_{\hat{\beta}} = 3$.
\item Find the MDE when $\sigma_{\hat{\beta}} = 1.5$.
\item Express the MDE as a general formula in terms of
$\sigma_{\hat{\beta}}$. What does this relationship imply about
how sample size affects the ability to detect small effects?
\end{enumerate}

\item 

\begin{enumerate}
\item Let $X$ be a continuous random variable with a strictly
increasing CDF~$F$. Show that $U = F(X)$ has a Uniform$(0,1)$ distribution.

\textit{Hints:} (i)~A random variable $U$ has a Uniform$(0,1)$
distribution if and only if $P(U \le u) = u$ for all
$u \in (0,1)$.
(ii)~Because $F$ is strictly increasing, its inverse $F^{-1}$
exists and satisfies $F(x) \le u \iff x \le F^{-1}(u)$ and
$F(F^{-1}(u)) = u$.

\item Let $Z \sim N(0,1)$ and define
$V = 2\bigl(1 - \Phi(|Z|)\bigr)$, where $\Phi$ is the standard
normal CDF.  Show that $V \sim \text{Uniform}(0,1)$.

\textit{Hints:} (i)~Start by computing $P(V \le v)$ for
$v \in (0,1)$.
(ii)~Note that $\Phi^{-1}(1 - v/2) = z_{1-v/2}$ is the critical
value of a two-sided test at significance level~$v$.

\item Consider testing $H_0\colon \beta = \beta_0$ against
$H_1\colon \beta \ne \beta_0$ with known variance
$\sigma_{\hat{\beta}}^2$.  Under $H_0$, the test statistic
$$
Z = \frac{\hat{\beta} - \beta_0}{\sigma_{\hat{\beta}}}
$$
has a standard normal distribution. The p-value of the two-sided
test is
$$
p = 2\bigl(1 - \Phi(|Z|)\bigr).
$$
Using Part~(b), show that the p-value has a Uniform$(0,1)$
distribution under $H_0$.

\item A test rejects $H_0$ whenever $p < \alpha$, where $\alpha \in (0,1)$. Using the resul from Part~(c),
show that the size of this test (i.e.,
$P(\text{reject } H_0 \mid H_0 \text{ is true})$) equals~$\alpha$. 
\end{enumerate}

\end{enumerate}
\end{document}
